\documentclass[12pt]{article}
\usepackage[utf8]{inputenc}
\usepackage[T1]{fontenc}
\usepackage[margin=1in]{geometry}
\usepackage{setspace}
\usepackage{xcolor}
\usepackage[colorlinks=true, urlcolor=blue!60!black, linkcolor=blue!60!black]{hyperref}
\setlength{\parindent}{2em}
\setlength{\parskip}{1.25em}
\renewcommand{\baselinestretch}{1.0}
\begin{document}
\begin{center}
{\LARGE \textit{Allen v.\ Cooper}, 589 U.S.\ 248 (2020)\par}
\vspace{0.8em}
{\large Case Brief\par}
\end{center}
\vspace{1em}

\section{Memory Jogger}
In a unanimous judgment arising from North Carolina's unauthorized online posting of shipwreck videos of Blackbeard's \textit{Queen Anne's Revenge}, the Supreme Court held that Congress lacked authority under either Article I's Intellectual Property Clause or Section~5 of the Fourteenth Amendment to abrogate state sovereign immunity in the Copyright Remedy Clarification Act of 1990 (\href{https://www.law.cornell.edu/supremecourt/text/18-877}{Cornell LII}; \href{https://supreme.justia.com/cases/federal/us/589/18-877/}{Justia}).

\section{Detailed Case Facts}
In 1717, the pirate Edward Teach---known as Blackbeard---captured a French slave ship in the West Indies and renamed her \textit{Queen Anne's Revenge}. In 1718, the vessel ran aground on a sandbar off Beaufort, North Carolina, where her remains lay undisturbed for nearly three centuries (\href{https://www.law.cornell.edu/supremecourt/text/18-877}{589 U.S.\ 248, Part~I}). In 1996, marine salvage company Intersal, Inc.\ discovered the shipwreck. Under the federal Abandoned Shipwreck Act, 43 U.S.C.\ \S~2105(c), and North Carolina law, N.C.\ Gen.\ Stat.\ Ann.\ \S~121-22, the wreck belonged to the State of North Carolina. The State contracted with Intersal to conduct recovery operations, and Intersal hired local videographer Frederick L.\ Allen (and his company, Nautilus Productions, LLC) to document the salvage. Over more than a decade, Allen produced videos and photographs of divers recovering cannons, anchors, and artifacts, and registered federal copyrights in all of those works (\href{https://www.law.cornell.edu/supremecourt/text/18-877}{Cornell LII}; \href{https://www.oyez.org/cases/2019/18-877}{Oyez}).

In 2013, Allen protested that North Carolina had uploaded his copyrighted works online without permission. The State settled that dispute for \$15,000 and executed an agreement defining the parties' respective rights. Shortly thereafter, however, Allen discovered that North Carolina had again posted five of his videos online and published one of his photographs in a state newsletter. In 2015, the North Carolina General Assembly also enacted N.C.\ Gen.\ Stat.\ \S~121-25(b)---dubbed ``Blackbeard's Law''---declaring photographs and video recordings of shipwrecks in the State's custody to be ``public record'' materials (\href{https://www.law.cornell.edu/supct/cert/18-877}{LII Bulletin}; \href{https://www.gwlr.org/allen-v-cooper-states-stay-copyright-pirates/}{Geo.\ Wash.\ L.\ Rev.\ On the Docket}). Allen sued Governor Roy A.\ Cooper, III and state officials in federal district court for copyright infringement and money damages.

\textbf{Key Constitutional and Statutory Provisions:}
\begin{itemize}
  \item \textbf{Intellectual Property Clause (U.S.\ Const.\ art.\ I, \S~8, cl.\ 8):} Empowers Congress ``[t]o promote the Progress of Science and useful Arts, by securing for limited Times to Authors and Inventors the exclusive Right to their respective Writings and Discoveries.''
  \item \textbf{Eleventh Amendment and Structural Sovereign Immunity:} The Eleventh Amendment bars federal judicial power over suits against a State by citizens of another State or foreign state, and the Court has long read it to confirm a broader structural presupposition shielding nonconsenting States from private federal suits, including by their own citizens (\href{https://supreme.justia.com/cases/federal/us/517/44/}{\textit{Seminole Tribe of Fla.\ v.\ Florida}, 517 U.S.\ 44 (1996)}).
  \item \textbf{Fourteenth Amendment, \S~1 (Due Process Clause) \& \S~5 (Enforcement Clause):} Prohibits States from depriving any person of ``life, liberty, or property, without due process of law,'' and grants Congress power ``to enforce, by appropriate legislation,'' the Amendment's substantive guarantees.
  \item \textbf{Copyright Remedy Clarification Act of 1990 (CRCA), 17 U.S.C.\ \S\S~501(a), 511(a)--(b):} Provides that a State ``shall not be immune, under the Eleventh Amendment [or] any other doctrine of sovereign immunity, from suit in Federal court'' for copyright infringement and exposes States to liability and remedies ``in the same manner and to the same extent as'' a private party (\href{https://www.law.cornell.edu/supremecourt/text/18-877}{Cornell LII}).
\end{itemize}

\textbf{Historical Context:} Congress enacted the CRCA in 1990 following a 1988 report by Register of Copyrights Ralph Oman (the ``Oman Report''), and two years later modeled the Patent Remedy Act on the CRCA's virtually identical text. In 1996, \href{https://supreme.justia.com/cases/federal/us/517/44/}{\textit{Seminole Tribe}} held that Article~I generally cannot abrogate state sovereign immunity, and in 1999, \href{https://supreme.justia.com/cases/federal/us/527/627/}{\textit{Florida Prepaid Postsecondary Ed.\ Expense Bd.\ v.\ College Savings Bank}, 527 U.S.\ 627 (1999)}, invalidated the Patent Remedy Act under both Article~I and Section~5. \textit{Allen} tested whether the CRCA could survive where its patent twin had fallen.

\section{Procedural History}
\begin{itemize}
  \item \textbf{U.S.\ District Court for the Eastern District of North Carolina:} Allen filed suit in December 2015 seeking damages for copyright infringement and challenging N.C.\ Gen.\ Stat.\ \S~121-25(b). North Carolina moved to dismiss on sovereign immunity grounds. The district court denied the motion (\textit{Allen v.\ Cooper}, 244 F.\ Supp.\ 3d 525 (E.D.N.C.\ 2017)), holding that the CRCA clearly expressed Congress's intent to abrogate and---although \href{https://supreme.justia.com/cases/federal/us/527/627/}{\textit{Florida Prepaid}} foreclosed Article~I abrogation---Congress validly abrogated immunity under Section~5 of the Fourteenth Amendment based on a ``pattern'' of abusive state copyright infringement (\href{https://www.law.cornell.edu/supremecourt/text/18-877}{Cornell LII}).
  \item \textbf{U.S.\ Court of Appeals for the Fourth Circuit:} On interlocutory appeal, the Fourth Circuit reversed (\textit{Allen v.\ Cooper}, 895 F.3d 337 (4th Cir.\ 2018)), holding that \href{https://supreme.justia.com/cases/federal/us/527/627/}{\textit{Florida Prepaid}} barred abrogation under Section~5 as well as Article~I because the CRCA's blanket abrogation was equally broad and rested on a similarly thin legislative record of unconstitutional infringement.
  \item \textbf{U.S.\ Supreme Court:} Because the Fourth Circuit invalidated a federal statute, the Supreme Court granted certiorari on June~3, 2019 (587 U.S.\ 980), heard oral argument on November~5, 2019, and \textbf{affirmed} the Fourth Circuit on March~23, 2020 (\href{https://www.oyez.org/cases/2019/18-877}{Oyez}). (Notably, after remand, the district court allowed Allen to reopen the case in 2021 on an as-applied theory under \textit{United States v.\ Georgia}, but in January 2026 the Fourth Circuit reversed and ordered dismissal with prejudice in \href{https://law.justia.com/cases/federal/appellate-courts/ca4/24-1954/24-1954-2026-01-23.html}{\textit{Allen v.\ Stein}, No.\ 24-1954 (4th Cir.\ 2026)}.)
\end{itemize}

\section{Judicial Votes}
\textbf{Judgment:} Unanimous (9--0) affirming the Fourth Circuit (\href{https://www.law.cornell.edu/supremecourt/text/18-877}{Cornell LII}).
\begin{itemize}
  \item \textbf{Majority Opinion:} Justice Kagan delivered the opinion of the Court, joined in full by Chief Justice Roberts and Justices Alito, Sotomayor, Gorsuch, and Kavanaugh (6 Justices), and joined by Justice Thomas except for the final paragraph in Part~II--A (stare decisis ``special justification'') and the final paragraph in Part~II--B (advice on future copyright legislation) (7 Justices for the core holding).
  \item \textbf{Concurring in Part and Concurring in the Judgment:} Justice Thomas filed an opinion concurring in part and concurring in the judgment.
  \item \textbf{Concurring in the Judgment:} Justice Breyer filed an opinion concurring in the judgment, joined by Justice Ginsburg.
  \item \textbf{Dissents:} None.
\end{itemize}

\section{Holding}
\textbf{Judgment:} Affirmed. North Carolina is entitled to state sovereign immunity from Allen's federal copyright infringement suit for money damages.

\textbf{Rule Announced:} Although Congress unequivocally expressed its intent to abrogate state sovereign immunity in the CRCA (17 U.S.C.\ \S\S~501(a), 511(a)), Congress lacked constitutional authority to do so. First, under \href{https://supreme.justia.com/cases/federal/us/517/44/}{\textit{Seminole Tribe}} and \href{https://supreme.justia.com/cases/federal/us/527/627/}{\textit{Florida Prepaid}}, Article~I's Intellectual Property Clause does not empower Congress to abrogate state sovereign immunity, and \href{https://supreme.justia.com/cases/federal/us/546/356/}{\textit{Central Va.\ Community College v.\ Katz}, 546 U.S.\ 356 (2006)}, is limited to the Bankruptcy Clause. Second, under Section~5 of the Fourteenth Amendment and \href{https://supreme.justia.com/cases/federal/us/521/507/}{\textit{City of Boerne v.\ Flores}, 521 U.S.\ 507 (1997)}, a blanket abrogation of immunity for all copyright infringement fails the ``congruence and proportionality'' test where the legislative record shows only a dozen possible examples of state infringement---mostly negligent---and fails to document a pattern of intentional or reckless infringement lacking adequate state-law remedies (\href{https://www.law.cornell.edu/supremecourt/text/18-877}{Cornell LII}).

\textbf{Holding vs.\ Dicta:}
\begin{itemize}
  \item \textbf{Holding:} The CRCA's across-the-board abrogation of state sovereign immunity is invalid under both Article~I, \S~8, cl.\ 8 and Section~5 of the Fourteenth Amendment; \textit{Florida Prepaid} controls and will not be overruled absent a ``special justification.''
  \item \textbf{Dicta:} (1) Justice Kagan's closing paragraph in Part~II--B inviting Congress to enact a future, tailored Section~5 copyright abrogation statute backed by a stronger legislative record to ``bring digital Blackbeards to justice''; (2) the characterization of \href{https://supreme.justia.com/cases/federal/us/546/356/}{\textit{Katz}} as a ``good-for-one-clause-only holding'' (subsequently undermined when the Court recognized additional Article~I ``plan of the Convention'' waivers in \textit{PennEast Pipeline Co.\ v.\ New Jersey} (2021) and \href{https://supreme.justia.com/cases/federal/us/597/20-603/}{\textit{Torres v.\ Texas Dept.\ of Public Safety}, 597 U.S.\ 580 (2022)}, as Justice Kagan herself acknowledged in her \textit{Torres} concurrence); and (3) the assumption---accepted by the parties without adversarial dispute---that copyrights qualify as ``property'' under the Fourteenth Amendment's Due Process Clause, which Justice Thomas expressly flagged as unresolved.
\end{itemize}

\section{Analysis of Opinions}

\subsection{Majority Opinion (Justice Kagan)}
Justice Kagan framed the case as controlled by precedent: ``\textit{Florida Prepaid} all but prewrote our decision today'' (\href{https://www.law.cornell.edu/supremecourt/text/18-877}{Cornell LII}). Applying the two-part abrogation framework, she first confirmed that the CRCA's text satisfied the clear-statement requirement, just as the verbatim Patent Remedy Act had in \href{https://supreme.justia.com/cases/federal/us/527/627/}{\textit{Florida Prepaid}}. She then rejected both constitutional bases asserted by Allen.

In \textbf{Part~II--A (Article~I)}, Justice Kagan explained that the Intellectual Property Clause covers patents and copyrights alike, and \href{https://supreme.justia.com/cases/federal/us/527/627/}{\textit{Florida Prepaid}}---building on \href{https://supreme.justia.com/cases/federal/us/517/44/}{\textit{Seminole Tribe}}---already held that Article~I cannot be used to circumvent sovereign immunity limits on federal jurisdiction. She rejected Allen's argument that \href{https://supreme.justia.com/cases/federal/us/546/356/}{\textit{Katz}} adopted a general ``clause-by-clause'' approach to Article~I abrogation. Instead, she read \textit{Katz} as reflecting ``bankruptcy exceptionalism'' rooted in bankruptcy's \textit{in rem} jurisdiction and unique Founding-era history, where the States surrendered immunity in the ``plan of the Convention'' itself without needing congressional abrogation. Moreover, because upholding the CRCA under Article~I would require overruling \textit{Florida Prepaid}, \textit{stare decisis} demanded a ``special justification'' over and above mere error, which Allen failed to offer.

In \textbf{Part~II--B (Fourteenth Amendment, Section~5)}, Justice Kagan applied \href{https://supreme.justia.com/cases/federal/us/521/507/}{\textit{City of Boerne}}'s ``congruence and proportionality'' test. While copyrights are property, a State violates the Due Process Clause only when infringement is intentional (or at least reckless) \textit{and} the State fails to provide an adequate state-law remedy (citing \textit{Daniels v.\ Williams} and \textit{Hudson v.\ Palmer}). Examining the 1988 Oman Report, she found the legislative record ``scarcely more impressive'' than in \textit{Florida Prepaid}: only twelve possible instances of state copyright infringement (seven lawsuits and five uncorroborated anecdotes), of which only two appeared intentional, coupled with testimony by Oman and congressional sponsors that state infringement was not widespread and mostly involved ``honest mistakes'' or ``sloppiness,'' and no inquiry into state-law remedies. Because the CRCA indiscriminately stripped immunity for every copyright claim, it was ``out of proportion'' to any documented due process violation. In closing, however, she noted that Congress could enact a tailored Section~5 statute in the future if it builds a record of unconstitutional state piracy.

\subsection{Concurring in Part and Concurring in the Judgment (Justice Thomas)}
Justice Thomas joined almost all of Justice Kagan's opinion but withheld his vote from two paragraphs and flagged one open constitutional question (\href{https://www.law.cornell.edu/supremecourt/text/18-877}{Cornell LII}). First, disagreeing with Part~II--A's final paragraph, he rejected the ``special justification'' standard for \textit{stare decisis}, reiterating his \textit{Gamble v.\ United States} view that Article~III obligates the Court to overrule any ``demonstrably erroneous'' precedent without extra hurdles---though \textit{Florida Prepaid} was correct on the merits, and he added in a footnote that \href{https://supreme.justia.com/cases/federal/us/546/356/}{\textit{Katz}} was ``not good for even that clause.'' Second, disagreeing with Part~II--B's final paragraph, he refused to offer advisory guidance to Congress on hypothetical future copyright statutes. Third, he emphasized that whether copyrights constitute ``property'' within the original meaning of the Fourteenth Amendment's Due Process Clause remains an open question, noting that \textit{Fox Film Corp.\ v.\ Doyal} addressed state taxation rather than due process.

\subsection{Concurring in the Judgment (Justice Breyer, joined by Justice Ginsburg)}
Justice Breyer concurred solely in the judgment on \textit{stare decisis} grounds (\href{https://www.law.cornell.edu/supremecourt/text/18-877}{Cornell LII}). Invoking \textit{Marbury v.\ Madison}, he underscored the anomaly of recognizing that States have a legal duty not to infringe copyrights while denying authors a federal damages remedy when States ``pirated intellectual property''---quipped with the example of a State hosting an unlicensed screening of \textit{Pirates of the Caribbean}. Reaffirming his long-held view that the Court ``went astray'' in \href{https://supreme.justia.com/cases/federal/us/517/44/}{\textit{Seminole Tribe}} and erred in \href{https://supreme.justia.com/cases/federal/us/527/627/}{\textit{Florida Prepaid}}, he nevertheless accepted that \textit{Florida Prepaid} controlled and expressed hope that Congress would venture into the ``great constitutional unknown'' to enact a tailored Section~5 remedy.

\section{Examples: Future Applications}
\begin{enumerate}
  \item \textbf{Same Side (Immunity Bars Suit):} A textbook publisher sues a public state university in federal court for money damages under the CRCA after a professor uploads scanned chapters to a university course portal. \textit{Why:} Just as in \textit{Allen}, the university is an arm of the State, the CRCA's blanket abrogation is invalid, and routine or negligent academic copying cannot trigger federal damages liability against the State.
  \item \textbf{Same Side (Immunity Bars Suit):} Congress passes a new federal statute abrogating state sovereign immunity for all federal trade-secret misappropriation claims under the Commerce Clause without building any legislative record of intentional state deprivations lacking state remedies. \textit{Why:} Under \textit{Allen}, \textit{Seminole Tribe}, and \textit{Florida Prepaid}, Article~I cannot support abrogation outside unique ``plan of the Convention'' exceptions, and an untailored statute without a Section~5 due process record fails \textit{City of Boerne}.
  \item \textbf{Opposite Side (Suit Permitted):} A photographer sues the director of a state tourism agency in her official capacity seeking strictly \textit{prospective injunctive relief} under \textit{Ex parte Young}, 209 U.S.\ 123 (1908), to stop the agency's ongoing unauthorized display of copyrighted photos on its website. \textit{Why:} \textit{Allen} bars retroactive money damages against the State under the CRCA, but does not disturb \textit{Ex parte Young} suits against state officers to enjoin ongoing violations of federal copyright law (\href{https://www.gwlr.org/allen-v-cooper-states-stay-copyright-pirates/}{Geo.\ Wash.\ L.\ Rev.\ On the Docket}).
  \item \textbf{Opposite Side (Suit Permitted):} A videographer sues a \textit{municipal} transit authority (a city agency, not an arm of the State) for copyright infringement damages in federal court. \textit{Why:} Eleventh Amendment and structural sovereign immunity protect States and arms of the State, not cities, counties, or independent local municipalities.
  \item \textbf{Fence-Sitter:} Congress enacts a narrow ``Willful State Copyright Piracy Act'' under Section~5, backed by a 2021 U.S.\ Copyright Office study documenting dozens of intentional state infringements and state-law immunity gaps, authorizing damages \textit{only} where a State engages in willful copyright infringement and provides no adequate state-court remedy. \textit{Why:} This directly tests Justice Kagan's closing roadmap in Part~II--B; while it remedies the CRCA's overbreadth, it remains a close question whether the Court would view the new empirical record as sufficient under \textit{Boerne} or whether originalist Justices would adopt Justice Thomas's reservation about whether copyrights are Fourteenth Amendment ``property.''
\end{enumerate}

\section{Critique}
\textbf{Scholarly Criticism and Logical Weaknesses:} Commentators across the ideological spectrum have criticized \textit{Allen}'s doctrinal foundations. \href{https://verdict.justia.com/2020/04/10/how-allen-v-cooper-breaks-important-new-if-dubious-ground-on-stare-decisis}{Professor Vikram David Amar} highlights a core analytical flaw inherited from \textit{Florida Prepaid}: the Court measures the \textit{overall frequency} of state IP infringement rather than the \textit{proportion} of covered conduct that is unconstitutional, even though a targeted cure should not require massive statewide abuse. Amar also critiques Justice Kagan's revival of \textit{Planned Parenthood v.\ Casey}'s ``special justification'' requirement for constitutional \textit{stare decisis}---citing a statutory precedent (\textit{Halliburton})---even where States have no legitimate reliance interest in pirating copyrights. Furthermore, \href{https://www.gwlr.org/allen-v-cooper-states-stay-copyright-pirates/}{Jasper L.\ Tran and Cameron Baker} note that stripping damages remedies guts the economic incentive to enforce copyrights against state actors, leaving creators with costly, fee-unrecoverable \textit{Ex parte Young} injunctions, while \href{https://digitalcommons.law.buffalo.edu/buffalolawreview/vol69/iss2/3/}{Buffalo Law Review scholarship} observes that \textit{Allen} deepens the bargaining-power imbalance between authors and state institutions. Finally, the Court's assertion that \href{https://supreme.justia.com/cases/federal/us/546/356/}{\textit{Katz}} was ``good-for-one-clause-only'' proved doctrinally unstable, collapsing just two years later in \href{https://supreme.justia.com/cases/federal/us/597/20-603/}{\textit{Torres v.\ Texas Dept.\ of Public Safety} (2022)}.

\textbf{Progressive / Pragmatic Perspective:} Progressive scholars and Justices Breyer and Ginsburg fault \textit{Allen} for entrenching the Rehnquist Court's ahistorical expansion of sovereign immunity in \href{https://supreme.justia.com/cases/federal/us/517/44/}{\textit{Seminole Tribe}} and \href{https://supreme.justia.com/cases/federal/us/527/627/}{\textit{Florida Prepaid}}. Under the Constitution's text, the Eleventh Amendment bars only diversity-style suits by out-of-state citizens, not federal-question suits enforcing exclusive federal copyrights that States themselves cannot adjudicate in state court (28 U.S.C.\ \S~1338(a)). By pairing federal-court exclusivity with state sovereign immunity, \textit{Allen} creates a right without a damages remedy, enabling ``state piracy'' in an era of digital public-university publishing (\href{https://www.everycrsreport.com/reports/LSB10465.html}{CRS Report LSB10465}).

\textbf{Originalist / Textualist Perspective:} Originalists are sharply divided. One camp---reflected in Justice Thomas's concurrence---supports \textit{Allen}'s protection of state sovereignty as a structural background principle (Hamilton's \textit{Federalist No.\ 81}) and applauds cabining \textit{Katz}, while criticizing the majority's heightened \textit{stare decisis} barrier and advisory dicta to Congress. Another textualist-originalist camp (e.g., scholars examining the Eleventh Amendment's text and Section~5's ``appropriate legislation'' language alongside \textit{McCulloch v.\ Maryland}) argues that neither \textit{Seminole Tribe}'s nontextual immunity bar nor \textit{Boerne}'s judge-made ``congruence and proportionality'' test tracks the original public meaning of the Eleventh or Fourteenth Amendments (\href{https://verdict.justia.com/2020/04/10/how-allen-v-cooper-breaks-important-new-if-dubious-ground-on-stare-decisis}{Amar, Justia Verdict}).

\section{Key Quotations}
\begin{quote}
``The slate on which we write today is anything but clean. \textit{Florida Prepaid}, along with other precedent, forecloses each of Allen's arguments.'' --- Justice Kagan, Opinion of the Court (\href{https://www.law.cornell.edu/supremecourt/text/18-877}{Cornell LII}).
\end{quote}

\begin{quote}
``But everything in \textit{Katz} is about and limited to the Bankruptcy Clause; the opinion reflects what might be called bankruptcy exceptionalism.'' --- Justice Kagan, Opinion of the Court (\href{https://www.law.cornell.edu/supremecourt/text/18-877}{Cornell LII}).
\end{quote}

\begin{quote}
``That kind of tailored statute can effectively stop States from behaving as copyright pirates. Even while respecting constitutional limits, it can bring digital Blackbeards to justice.'' --- Justice Kagan, Opinion of the Court (\href{https://www.law.cornell.edu/supremecourt/text/18-877}{Cornell LII}).
\end{quote}

\begin{quote}
``If our decision in \textit{Florida Prepaid} were demonstrably erroneous, the Court would be obligated to `correct the error, regardless of whether other factors support overruling the precedent.'\,'' --- Justice Thomas, concurring in part and concurring in the judgment (\href{https://www.law.cornell.edu/supremecourt/text/18-877}{Cornell LII}).
\end{quote}

\begin{quote}
``In its view, Congress' power under the Intellectual Property Clause cannot support a federal law providing that, when proven to have pirated intellectual property, States must pay for what they plundered.'' --- Justice Breyer, concurring in the judgment (\href{https://www.law.cornell.edu/supremecourt/text/18-877}{Cornell LII}).
\end{quote}

\end{document}
